% Primitive: independence — authored from probability/independence_manim.py
% and its README concepts table (six scenes, watched in order), plus plan
% 002's verified anchors. Backward-grounded: the counting and alignment
% chapters only; the conditional lens belongs to the next chapter. Problem
% answers are computed by answers/independence.py (the solve-gate anchor);
% new anchor keys ride primitives/anchors-proposals-probability.yaml until
% merged into anchors.yaml.
\chapter{Independence}

\begin{narrative}
The counting chapter multiplied option counts and owed nothing for it.
The alignment chapter multiplied per-frame scores into a path's weight —
and flagged the move as an assumption it could not yet defend. This
chapter is where multiplying becomes lawful for probabilities: what the
product of two probabilities actually asserts, how to check it, the two
confusions it attracts, and what assuming it costs.

\section{Probability is area}

Draw the sample space as a $1 \times 1$ square: an \emph{event} is a
region, and its probability is its share of the
area~\cite{probability-3blue1brown-bayes-theorem-lesson}. A fair die is
six equal cells; the event ``even'' covers three of them, and
$3/6$ — count over total — is the same number as the region's area.
Nothing new was computed: proportion is what counting was already doing,
divided through by the size of the whole. The math of probability is the
math of proportions, and every argument in this chapter is an argument
about regions of this square.

\section{The product rule}

Run two experiments — a coin and a die, two dice, two frames. The
counting chapter's grid returns with its cells reweighted from counts to
areas: an event about the \emph{first} experiment selects rows, an event
about the \emph{second} selects columns, and their intersection is a
rectangle. Rectangle area is width times height — that is the whole
proof. On the unit square, both cuts run straight across, all the
way~\cite{probability-tsvibt-jaime-sevilla-mollina-two-independent-events-square-v}.

\begin{center}
\begin{tikzpicture}[scale=0.9]
  % The aligned unit square — both cuts straight across (scene 2's
  % device, carrying scene 3's jewel numbers). B-band Cool, the
  % intersection rectangle Accent: the result being built toward.
  \fill[Cool!12] (0,0) rectangle (4,2);
  \fill[Cool!22] (0,0) rectangle (8/3,4);
  \fill[AccentInk!30] (0,0) rectangle (8/3,2);
  \draw[Muted] (0,0) rectangle (4,4);
  \draw[Muted] (8/3,0) -- (8/3,4);
  \draw[Muted] (0,2) -- (4,2);
  \node[below, color=CoolInk] at (4/3,-0.4) {$P(B)=\tfrac23$};
  \node[left, color=CoolInk] at (-0.45,1) {\rotatebox{90}{$P(A)=\tfrac12$}};
  \node[right, color=AccentInk] at (4.45,1)
    {$P(A\cap B)=\tfrac12\cdot\tfrac23=\tfrac13$};
\end{tikzpicture}
\end{center}

When the picture looks like this — each event a straight band, the
overlap a rectangle — the joint probability factors, and that factoring
is the \emph{definition} of
independence~\cite{probability-wikipedia-independence-probability-theory}:
\[
  A \perp B
  \quad\Longleftrightarrow\quad
  P(A \cap B) \;=\; P(A)\,P(B).
\]
As in the counting chapter, the formula arrives last: first the grid,
then the rectangle, then the equation that summarises what the picture
already said. The definition is the product form itself — not a
statement about causes, mechanisms, or time.

\section{One die, two events}

Independence does not need two experiments. On a single fair die, take
$A = \{\text{even}\}$ and $B = \{1,2,3,4\}$. Then
$A \cap B = \{2,4\}$, and
\[
  P(A \cap B) = \tfrac{2}{6} = \tfrac13
  = \tfrac12 \cdot \tfrac23 = P(A)\,P(B)
\]
— independent, though both events are about the same roll. Now slide
$B$'s boundary a single pip, to $\{1,2,3\}$: the intersection shrinks to
$\{2\}$, and $\tfrac16 \neq \tfrac12 \cdot \tfrac12 = \tfrac14$.
Dependent. One pip decided it. And the same two events that factored on
the fair die fail to factor on a biased one: independence is a property
of the \emph{measure} — of how the weight is spread — not of the events
alone~\cite{probability-kevin-davisross-probability-and-simulation-3-5-independence}.
``Separate mechanisms'' is neither necessary for independence nor
sufficient; the arithmetic is the only test.

\section{Mutually exclusive is not independent}

The most-tested confusion in probability: surely events that cannot
co-occur are ``unrelated''? The product test says the opposite, as
loudly as it can. If $A$ and $B$ are disjoint with positive
probability, then
\[
  P(A \cap B) \;=\; 0 \;\neq\; P(A)\,P(B) \;>\; 0 .
\]
Disjoint regions share no area, so the left side is zero while the
right side is not — mutually exclusive events are \emph{maximally}
dependent: seeing one tells you the other did not happen. On the die,
$\{\text{even}\}$ and $\{1,3,5\}$ fail the test at $0$ versus
$\tfrac14$. (One honest exception: an event of probability zero is
trivially independent of everything.) This is also why Venn-style
pictures mislead here — disjoint circles \emph{look} unrelated. The
square tells the truth the circles hide: dependence is the cut through
the band stepping; independence is the knife-edge case where it runs
straight.

\section{Chains of trials}

One flip cuts the square in half. A second flip cuts each half in half
again; after four flips the square holds sixteen cells, and the
particular sequence HHTH is one cell of area
$\anchor{002.hhth}$ — each independent trial multiplied one more factor
on. In general, for independent events $A_1, \dots, A_n$,
\[
  P(A_1 \cap \cdots \cap A_n) \;=\; \prod_{i=1}^{n} P(A_i),
\]
a box with side lengths $p_1, \dots, p_n$. Two honest captions come
with the chain. First, it needs \emph{mutual} independence, which is
strictly more than every pair factoring: with two fair coins, take
$A = \{\text{first is H}\}$, $B = \{\text{second is H}\}$,
$C = \{\text{exactly one head}\}$ — all three pairs factor at
$\tfrac14$ each, yet the triple intersection is empty, $0 \neq \tfrac18$
(Bernstein's construction~\cite{probability-wikipedia-pairwise-independence}).
Second, asserting the product over a real sequence of trials is a
\emph{modeling choice} — the product
measure~\cite{probability-mit-6-041-lecture-13-the-bernoulli-process} —
not an observation. This is precisely the purchase the alignment
chapter made on credit when it priced a path as a product of per-frame
scores: one factor per frame is this formula, asserted about a model.

\section{When to use it}

The chapter's exit is a field guide. \emph{Replacement installs
independence}: draw a card, put it back, shuffle — the second draw
faces the same deck, and two aces cost
$\anchor{002.aces.replaced}$. \emph{Depletion breaks it}: keep the
first card out and the same two aces cost
$\anchor{002.aces.depleted}$ — the counting chapter's shrinking pool,
returned exactly as promised, as the thing that breaks independence.
\emph{Common causes break it too}, with no causal link between the
events themselves: two symptoms of one disease co-occur far more often
than their product predicts. The alignment chapter already met this in
the wild — within a word, the frames of speech share the word as a
common cause, which is why real systems bolt an external language model
onto a per-frame emitter at decode time. And the \emph{gambler's
fallacy} is independence misread in reverse: a run of heads owes the
future nothing; the long-run proportion settles by \emph{swamping} —
ever more trials diluting any early excess — not by
compensation~\cite{probability-libretexts-the-gambler-s-fallacy}.
Independence is a property of the model's measure, and assuming it is a
purchase. Know what you paid.

\paragraph{Where this goes.} The next chapter buys back this chapter's
shortcuts: the stepped cut gets a name — conditional probability — the
aces' $\anchor{002.aces.depleted}$ finally gets its license as a
product of a probability and a \emph{conditional} probability, and the
assumption the alignment chapter actually needs is stated in full:
frames independent \emph{given the input}. Further down the road, the
sixteen-cell square built here is stamped and sorted into a
distribution, and the per-frame scores become honest probabilities.
\end{narrative}

\section*{Practice problems}

\begin{problem}
Roll two fair dice. Let $A = \{\text{first die shows 5 or 6}\}$ and
$B = \{\text{second die shows 3 or less}\}$. Compute $P(A \cap B)$ and
check it against the product rule.
\begin{hint}
Draw the $6 \times 6$ grid; $A$ selects rows, $B$ selects columns.
\end{hint}
\begin{solution}
The intersection is a rectangle of $2 \times 3 = 6$ cells out of $36$:
$P(A \cap B) = \tfrac{6}{36} = \tfrac16$. The margins give
$P(A) = \tfrac13$ and $P(B) = \tfrac12$, and
$\tfrac13 \cdot \tfrac12 = \tfrac16$ — the product rule holds, because
an event about the first die and an event about the second always cut
the grid as perpendicular bands, and a rectangle's share of the grid is
width times height.
\end{solution}
\end{problem}

\begin{problem}
On one fair die, test $A = \{\text{even}\}$ for independence against
$B = \{1,2,3,4\}$ and against $B' = \{1,2,3\}$.
\begin{hint}
Three probabilities each time; the definition is the only test.
\end{hint}
\begin{solution}
Against $B$: $A \cap B = \{2,4\}$, so
$P(A \cap B) = \tfrac13 = \tfrac12 \cdot \tfrac23$ — independent,
inside a single roll. Against $B'$: $A \cap B' = \{2\}$, so
$P(A \cap B') = \tfrac16$ while
$P(A)P(B') = \tfrac12 \cdot \tfrac12 = \tfrac14$ — dependent. Moving
one pip from $B$ to $B'$ flipped the verdict: independence is
arithmetic about the measure, not intuition about mechanisms.
\end{solution}
\end{problem}

\begin{problem}
Show that $A = \{\text{even}\}$ and $B = \{1,3,5\}$ on a fair die are
not independent, and say how much information seeing $B$ carries about
$A$.
\begin{hint}
What is the area of a disjoint overlap?
\end{hint}
\begin{solution}
$A \cap B = \emptyset$, so $P(A \cap B) = 0$, while
$P(A)P(B) = \tfrac12 \cdot \tfrac12 = \tfrac14$. The product test fails
as loudly as it can: knowing $B$ occurred settles $A$ completely (it
did not happen). Mutually exclusive events with positive probability
are never independent — they are maximally dependent.
\end{solution}
\end{problem}

\begin{problem}
Flip two fair coins. Let $A = \{\text{first is H}\}$,
$B = \{\text{second is H}\}$, $C = \{\text{exactly one head}\}$. Verify
that every pair is independent, and that the three events are not
mutually independent.
\begin{hint}
Four outcomes, each of probability $\tfrac14$; check the three pair
products, then the triple.
\end{hint}
\begin{solution}
$P(A) = P(B) = P(C) = \tfrac12$. Pairs: $A \cap B = \{\mathrm{HH}\}$,
$A \cap C = \{\mathrm{HT}\}$, $B \cap C = \{\mathrm{TH}\}$ — each of
probability $\tfrac14 = \tfrac12 \cdot \tfrac12$, so all three pairs
factor. But $A \cap B \cap C$ demands two heads \emph{and} exactly one
head: impossible, so its probability is $0 \neq \tfrac18$. Pairwise
independence does not give mutual independence — which is why the chain
formula's condition says \emph{mutual}.
\end{solution}
\end{problem}

\begin{problem}[stretch]
Draw two cards from a shuffled deck. Compute $P(\text{both aces})$ with
replacement and without, and — without replacement — $P(\text{second
card is an ace})$ before anyone looks at the first.
\begin{hint}
Replacement makes the draws independent trials; depletion is the
shrinking pool. For the last part, enumerate — or argue by symmetry.
\end{hint}
\begin{solution}
With replacement the draws are independent trials:
$\left(\tfrac{1}{13}\right)^2 = \anchor{002.aces.replaced}$. Without,
the pool shrinks: $\tfrac{4}{52} \cdot \tfrac{3}{51} =
\anchor{002.aces.depleted}$ — smaller, because the first ace depletes
the aces. Yet $P(\text{second is ace}) = \tfrac{1}{13}$ exactly
($204/2652$ enumerated): as long as no one looks at the first card, the
second is just a uniformly random card. Depletion changes the
\emph{joint} behaviour, not the margin — which is exactly what
``dependent'' means.
\end{solution}
\end{problem}
